% NOTE:
% latexmk -shell-escape -pvc poster.tex # Watches and compiles on each change.
% latexmk -c poster.tex   # Clean the temporal files.

\documentclass[20pt]{beamer}

% Adapted from;
% How to Design a Scientific Poster using Beamer -1 (Latex Basic Tutorial-30)
% https://youtu.be/2ZWnFFhVkdE?si=N3e1Ob8YWyTwW19e

\usepackage[size=custom, width=90, height=120, orientation=portrait, scale=1.4]{beamerposter}
\usetheme{Madrid}
\usepackage{changepage}
\usepackage[numbers]{natbib}
\usepackage{listings}
\usepackage{siunitx}
\usepackage{hyperref}
\usepackage[linesnumbered,algoruled,longend]{algorithm2e}
\usepackage{graphicx}
\usepackage{subcaption}
\usepackage{lipsum}

\usepackage{amsmath}
\usefonttheme[onlymath]{serif}



\beamertemplatenavigationsymbolsempty
\newcommand{\R}{R}
\renewcommand{\raggedright}{\leftskip=0pt \rightskip=0pt}
\let\olditemize\itemize
\renewcommand{\itemize}{\olditemize\addtolength{\itemsep}{0.4\baselineskip}}

\addtobeamertemplate{block begin}{}{\vspace{5mm}\begin{adjustwidth}{5mm}{5mm}}
\addtobeamertemplate{block end}{\end{adjustwidth}\vspace{10mm}}{\vspace{5mm}}

\setbeamertemplate{caption}[numbered]

\setbeamerfont{block title}{size={\centering\bfseries\fontsize{48}{60}}}
\setbeamercolor{block body}{bg=white}
\setbeamercolor{background canvas}{bg=white}
%\definecolor{title-color}{RGB}{23, 80, 22}
%\setbeamercolor{fbcolor}{fg=white, bg=cyan}

\definecolor{titlebg}{RGB}{23,80,20}
\definecolor{rulebg}{RGB}{248,230,0}
\definecolor{titlefontcolor}{RGB}{239,222,135}

\setbeamercolor{block title}{bg=titlebg,fg=titlefontcolor}

\setbeamercolor{fbcolor}{fg=white, bg=titlebg}
\setbeamertemplate{headline}{
    \begin{beamercolorbox}[wd=\paperwidth]{fbcolor}\vskip5mm
        \begin{columns}

            \begin{column}{0.20\linewidth}
                \vspace{5mm}
            \centering
            \includegraphics[scale=0.7]{logos/logo_tum.png}\\
                \vspace{2cm}
                \includegraphics[scale=0.6]{logos/logo_baylat.png}
            \end{column}

            \begin{column}{0.60\linewidth}
                \vspace{5mm}
            \centering\bfseries
\textcolor{titlefontcolor}{
\fontsize{80pt}{96} \selectfont Estimation of the tropical fire season\\ [6mm]
{\fontsize{48pt}{60}\selectfont 
\underline{Alber Sanchez}${}^1$, Guilherme Mataveli${}^1$, \\
                Luiz E. O. C. Arag\~{a}o${}^1$ 
\\}
\vspace{5mm}
{\fontsize{38}{48}\selectfont ${}^1$National Institute for Space Research.\\} 
\vspace{5mm}
{\fontsize{44}{54}\selectfont alber.ipia@inpe.br}
                }
            \end{column}

            \begin{column}{0.20\linewidth}
                \vspace{5mm}
            \centering
            \includegraphics[scale=2.0]{logos/logoinpe-azul-menor.png} \\
               \vspace{2cm}
                \includegraphics[scale=2.0]{logos/logo_fapesp.png}
               \hspace{5mm}
            \end{column}

        \end{columns}\vskip 5mm
    \end{beamercolorbox}
    \color{rulebg}\rule{\paperwidth}{20pt}
}

\setbeamercolor{footcolor}{fg=white, bg=black}
\setbeamertemplate{footline}{
    \begin{beamercolorbox}[wd=\paperwidth, center, ht=1.5em]{footcolor}
        \small Any questions or suggestions? Write to alber.ipia@inpe.br
    \end{beamercolorbox}
}

\begin{document}\vspace*{-2cm}
\begin{frame}[fragile,t]
\begin{columns}[t]


%==== Column 1 ====

\begin{column}{0.33\linewidth}

    \vspace{1.0cm}

    \begin{block}{Introduction\vphantom{g}}

    \vspace{1.0cm}

Wildfires in tropical rainforests degrade, among other essential functions, their carbon storage capacity, which plays a critical role in mitigating the effects of the ongoing climate crisis.
Improved forecasting of the fire seasons in these forests would provide better data for decision-makers, supporting the development of adaptation policies to address issues such as excessive wildfire smoke and declining rainfall and drinking water availability.
This study presents a pixel-wise estimation of the fire season in tropical forests, utilizing VIIRS fire hotspots and sigmoidal models.

    \end{block}

    \vspace{0.5cm}

    \begin{block}{Data\vphantom{g}}

\vspace{1.0cm}

\begin{itemize}
    \item The near real-time Suomi-NPP Visible Infrared Imaging Radiometer Suite (VIIRS) Active Fire Detection product~\cite{nasaviirslandscienceteam2020}, processed by the University of Maryland with a three-month lag and distributed through the Fire Information for Resource Management System, provides active fire detection data at a nominal spatial resolution of 375 meters.
    \item The high resolution K\"{o}ppen-Geiger maps based on constrained CMIP6 projections~\cite{beck2023}.
    \item The R package \textit{sicegar} for fitting double-sigmoidal models to the data~\cite{caglar2018}.
\end{itemize}

    \end{block}

    \vspace{0.5cm}

    \begin{block}{Double Sigmoidal model\vphantom{g}}

\vspace{1.0cm}

It is based on the logistic function (Equation~\ref{eq:logistic_function}) where $I(t)$ is the intensity as a function of time; $I_{max}$ and $t_{mid}$ are the maximum observed intensity and the time when it is reached; $a_{1}$ is the slope of $I(t)$ at $t_{mid}$ given by the formula $\frac{d}{dt} I(t)t = t_{mid} = a_{1}I_{max}/4$.

\begin{equation}
    I(t) = f_{sig}(t) = \frac{I_{max}}{ 1 + exp(-a_{1}(t - t_{mid}))}
    \label{eq:logistic_function}
\end{equation}

    \vspace{0.5cm}

The double-sigmoidal is obtained by multiplying two sigmoidal functions (Equation~\ref{eq:two_sigmoidal}).

\begin{equation}
    {
        \scriptstyle{
            f_{dsig-base}(t) = \frac{1}{1 + exp(-a'_{1}(t -t'_{mid1}))}
            \frac{1}{1 + exp(-a'_{2}(t -t'_{mid2}))}
        }
        \label{eq:two_sigmoidal}
    }
\end{equation}

    \vspace{0.5cm}

And then let $t^*$ be the time at which $f_{dsig-base}(t)$ is maximal and $f_{max} = f_{dsig-base}(t^*)$.
In this way we can write:

\begin{equation*}
    {
        \scriptstyle{
            I(t) = f_{dsig}(t) = 
            \left\{
                \begin{array}{rl}
                    c_{1}  f_{dsig-base}(t) \text{ for } t <= t^* \text{(growth phase)} \\
                    c_{2}  f_{dsig-base}(t) + I_{final} \text{ for } t > t^* \text{(decay phase)}
                \end{array}
            \right.
        }
        }
\end{equation*}

    \vspace{0.5cm}

Where $c_{1} = \frac{I_{max}}{f_{max}}$ and $c_{2} = \frac{(I_{max} - I_{final})}{f_{max}}$.

    \vspace{0.5cm}

The features of this model are displayed in Figure~\ref{fig:double-sigmoidal_model}.

\begin{figure}[ht]
    \centering
    \includegraphics[width=0.9\textwidth]{images/caglar2018_double_sigmoidal_model.png}
    \caption{Double-sigmoidal model. Their midpoints are assumed as the start and end of the fire season. Source:~\cite{caglar2018}.}
    \label{fig:double-sigmoidal_model}
\end{figure}

    %\vspace{0.5cm}

    \end{block}

\end{column}



%==== Column 2 ====

\begin{column}{0.33\linewidth}

    \vspace{1.0cm}

    \begin{block}{Method\vphantom{g}}
\vspace{1cm}
\begin{itemize}
    \item We aggregated VIIRS Active Fire data (from 2012 to 2020) using a regular grid made of 0.25-degree-month cells over a WGS84 Coordinate Reference System.
    \item Then we estimated, for each month and pixel, the mean over 9 years of data.
    \item Then we fitted a double-sigmoidal function, assuming that their mid-points are the start and end of the fire season and their durations are the time intervals between midpoints. 
    \item Finally, our results were cropped using the K\"{o}ppen-Geiger tropical monsoon and rainforest zones for the period 1991-2020~\cite{beck2023}.
\end{itemize}

    \end{block}

    \vspace{0.5cm}

    \begin{block}{Results\vphantom{g}}

    \vspace{0.5cm}

Our results show that the fire season in the tropical forests lasts between one and three months.
Additionally, both the start and end of the season displays a horizontal striping pattern, depending on the latitude and relative position to the equatorial line (Figures~\ref{fig:dsig_tropical_africa}, \ref{fig:dsig_tropical_america}, and \ref{fig:dsig_tropical_asia}).

    \vspace{0.5cm}

\begin{figure}[ht]
\centering
\includegraphics[width=0.9\textwidth]{figures/dsig_tropical_africa.png}
\caption{Fire season in the African tropical rainforest.}
\label{fig:dsig_tropical_africa}
\end{figure}

    \vspace{0.5cm}

\begin{figure}[ht]
\centering
\includegraphics[width=0.9\textwidth]{figures/dsig_tropical_america.png}
\caption{Fire season in the American tropical rainforest.}
\label{fig:dsig_tropical_america}
\end{figure}

    \vspace{0.5cm}

\begin{figure}[ht]
\centering
\includegraphics[width=0.9\textwidth]{figures/dsig_tropical_asia.png}
\caption{Fire season in the tropical rainforest of Asia and Oceania.}
\label{fig:dsig_tropical_asia}
\end{figure}

    \vspace{0.5cm}

Figure~\ref{fig:plot_latitude} (left) clearly shows the differences between the Northern and Southern Hemispheres regarding the fire season start.
This pattern has already been observed in the Amazon region using climatological variables rather than Active Fire data~\cite{carvalho2021a}.

    \vspace{0.5cm}

\begin{figure}[ht]
\centering
\includegraphics[width=0.45\textwidth]{figures/plot_latitude_pos_from.png}
\includegraphics[width=0.45\textwidth]{figures/plot_latitude_val_len.png}
\caption{Start and duration of the fire season aggregated by latitude.}
\label{fig:plot_latitude}
\end{figure}

    %\vspace{0.5cm}

    \end{block}

\end{column}



%==== Column 3 ====

\begin{column}{0.33\linewidth}

    \vspace{1cm}

    \begin{block}{Reproducibility\vphantom{g}}
        \vspace{0.5cm}
        To ensure the reproducibility of our results, we used only publicly available data and software.
        Besides the previously introduced datasets, we used the \R{} programming language for statistical computing and data visualization~\cite{ihaka1996}.
        We also used the \R{} packages \textit{sf} and \textit{terra} for handling vector and raster data and the \textit{sicegar} for fitting double-sigmoidal models~\cite{pebesma2018,hijmans2020,caglar2018}.
    \end{block}

    \begin{block}{Future work\vphantom{g}}

        \vspace{0.5cm}

We view our results as a starting point to addressing the question: \textit{Is the fire season changing?}
To find an answer, we will:
\begin{itemize}
    \item Switch from VIIRS to MODIS active fire data to get longer time series.
    \item Run the analysis presented here on a time series longer than 20 years and use it as a baseline.
    \item Compute yearly deviations from the baseline.
    \item Test for patterns in the deviations, particularly changes in the start and duration of the fire season.
    \item Compare our estimation of the fire season to the dry season. This could point to human-behavior changes since most of the rainforest fires are anthropogenic.
\end{itemize}
    \end{block}

\vspace{0.5cm}

    \begin{block}{References\vphantom{g}}

\vspace{0.5cm}    

\bibliography{workshop_tropical_vegetation}
\bibliographystyle{sbc}
    \end{block}

\vspace{0.5cm}

    \begin{block}{Links}

\vspace{0.5cm}

\begin{figure}
    \begin{subfigure}[b]{0.2\textwidth}
\centering
\includegraphics[width=0.90\textwidth]{images/qrcode_email_alber_ipia_at_inpe.png}\\
{Email.}
    \end{subfigure}
    ~
    \begin{subfigure}[b]{0.2\textwidth}
\centering
        \includegraphics[width=0.88\textwidth]{images/qrcode_seasonmetrics.png}
{Code.}
    \end{subfigure}
\end{figure}

% \begin{itemize}
% \item Email: \href{mailto:alber.ipia@inpe.br}{alber.ipia@inpe.br} 
% \item Code: {\scriptsize\url{https://github.com/albhasan/seasonmetrics}}
% \end{itemize}

    \end{block}
\end{column}

\end{columns}
\end{frame}
\end{document}
